% c01 —— §1 末 + §2 K3 曲面上的镜像对称 + §3 Batyrev 镜像构造与 Dolgachev 猜想
%        （p002–p003 / 原书 p.128–129）
% 承 c00 末句「……之间的奇妙关系。」；原书 p.127 的末句「これは Gromov-Witten 不変量を」
% 在页边界截断，整句归本块，从「这是把 Gromov–Witten 不变量用于……」起写。

% ===== 原书 p.128 =====

这是把 Gromov–Witten 不变量用于在 $Y$ 的「Kähler 结构的模空间」上确定的 Hodge 结构变形
（$A$-VHS），与在 $\check{Y}$ 的复结构模空间上按通常方法确定的 Hodge 结构变形（$B$-VHS）
之间的关系来表述，因此有时也称为 \textbf{Hodge 理论镜像对称}（例如参见 [CK99, Iri11] 及其参考文献）。
这里给「Kähler 结构的模空间」加上引号，是因为目前它的定义还没有确定。人们预期，
$Y$ 的「Kähler 结构的扩大模空间」的万有覆盖，会成为 $Y$ 上相干层的导出范畴的
Bridgeland 意义下的稳定性条件的空间。这里所说的「扩大模空间」，是指不仅考虑
$h^{1,1}(Y)$，而且考虑全部代数闭链。\textbf{同调镜像对称}是 Kontsevich 在 1994 年
国际数学家大会上提出的猜想，它断言三角范畴的等价
\begin{equation}
  D^{b}\mathfrak{F}\mathrm{uk}\,Y \cong D^{b}\operatorname{coh}\check{Y} \tag{1.2}
\end{equation}
的存在。\textbf{Strominger–Yau–Zaslow 猜想}断言：任意 Calabi–Yau 流形 $Y$ 都具有特殊
Lagrange 环面纤维丛 $Y\to B$ 的结构，并把对偶环面纤维丛 $\check{Y}\to B$ 的全空间
作为镜像 $\check{Y}$ 得到。就难度而言，一般认为
\begin{center}
（拓扑镜像对称）$\ll$（古典镜像对称）\\
$\ll$（同调镜像对称）$\ll$（SYZ 猜想）
\end{center}
的程度，但对给定空间构造其镜像一般都很难，就连拓扑镜像对称也远未完全解决。
至于镜像的构造，已知有
\begin{itemize}[leftmargin=2.2em]
  \item 轨形构造 [GP90]
  \item 自反多面体的极对偶 [Bat94]
  \item 可逆多项式的转置 [BH93]
  \item 环面退化与 conifold 转迁 [BCFKvS98]
  \item 特殊 Lagrange 环面纤维丛 [SYZ96]
\end{itemize}
等各种方法。

\section{K3 曲面上的镜像对称}

设 $Y$ 为 K3 曲面。在自由阿贝尔群 $H^{*}(Y;\mathbb{Z})$ 上赋予向井配对
\[
  ((a_{0},a_{2},a_{4}),(b_{0},b_{2},b_{4}))=(a_{2},b_{2})-(a_{0},b_{4})-(a_{4},b_{0})
\]
后得到的格称为扩大 K3 格。$Y$ 的复结构由满足 Hodge–Riemann 双线性关系式
\begin{equation}
  (\Omega,\Omega)=0,\qquad(\Omega,\overline{\Omega})>0 \tag{2.1}
\end{equation}
的全纯体积形式 $\Omega\in H^{2,0}(Y)$ 的类 $[\Omega]\in\mathbb{P}(H^{2}(Y;\mathbb{C}))$ 确定。

% ===== 原书 p.129 =====

所谓 $Y$ 的复化 Kähler 结构，是指 $H^{*}(Y;\mathbb{C})$ 中的元素
\begin{align*}
  \mho&=\exp(B+\sqrt{-1}\omega)\\
  &=\Bigl(1,\ B+\sqrt{-1}\omega,\ \frac{1}{2}(B+\sqrt{-1}\omega)^{2}\Bigr)\\
  &\in H^{0}(Y;\mathbb{C})\oplus H^{2}(Y;\mathbb{C})\oplus H^{4}(Y;\mathbb{C})
\end{align*}
且满足
\begin{equation}
  (\mho,\mho)=0,\qquad(\mho,\overline{\mho})>0,\qquad
  (\mho,\Omega)=0,\qquad(\mho,\overline{\Omega})=0 \tag{2.2}
\end{equation}
者。$H^{2}(Y;\mathbb{R})$ 中的元素 $\omega$ 是 Kähler 类的推广，满足
$\omega\in H^{2}(Y,\mathbb{R})\cap H^{1,1}(Y)$ 以及 $\omega\cdot\omega>0$。
$H^{2}(Y,\mathbb{R})$ 中的元素 $B$ 称为 B 场，满足
$B\in H^{2}(Y,\mathbb{R})\cap H^{1,1}(Y)$。注意 $\omega$ 与 $B$ 由
$[\mho]\in\mathbb{P}(H^{*}(Y;\mathbb{C}))$ 唯一确定。把 K3 曲面 $Y$ 与
$\mathbb{P}(H^{*}(Y;\mathbb{C}))$ 的两个元素所成的组 $(Y,([\Omega],[\mho]))$
且满足上述条件者，称为带有复化 Kähler 结构的 K3 曲面。

\begin{definition}[Aspinwall–Morrison {[AM97]}]
带有复化 Kähler 结构的 K3 曲面的组
\[
  \bigl((Y,([\Omega],[\mho])),\ (\check{Y},([\check{\Omega}],[\check{\mho}]))\bigr)
\]
若存在扩大 K3 格的同构
$\varphi\colon H^{*}(Y;\mathbb{Z})\xrightarrow{\sim}H^{*}(\check{Y};\mathbb{Z})$ 使得
$(\varphi([\Omega]),\varphi([\mho]))=([\check{\mho}],[\check{\Omega}])$，
则称其为 \textbf{Aspinwall–Morrison 意义下的镜像对}。
\end{definition}

\section{Batyrev 镜像构造与 Dolgachev 猜想}

设 $X$ 为光滑环面弱 Fano 流形。这里，弱 Fano 流形是指反典范除子为 nef 且巨大的流形。
设 $Y$ 为 $X$ 的光滑反典范超曲面。反典范丛的 Newton 多面体定义为
\begin{equation}
  \Delta=\Conv\bigl\{m\in M\,\bigm|\,x^{m}\in H^{0}(\mathcal{O}_{X}(-K_{X}))\bigr\} \tag{3.1}
\end{equation}
这里 $M$ 是 $X$ 的稠密环面的特征标群：
\begin{equation}
  \mathbb{T}=\Spec\mathbb{C}[M]\subset X. \tag{3.2}
\end{equation}

% CHUNK c01 DONE pages p002–p003
